\section{An exact scale flow and the temperature curve}\label{sec:flow}

Theorem~\ref{thm:transfer} controls the linear effect of moving a bounded
potential to a larger square.  We now implement that operation by an exact
change of the bulk interaction.  Its nonlinear error will be quadratic,
and its only expanding direction will be removed by choosing the temperature.
After estimating the residual caused by periodic evaluation, reversing the
flow will couple the resulting torus law to the pure law by edits in
independently marked squares.

The bulk scale parameters will use only the pure full-plane law
\(\mu _0\), fixed scale ratios, and the given interaction \(U\),
independently of the target domain and its boundary data.  All square
coordinates are lattice coordinates.  We write
\[
 \begin{gathered}
 B_r(x)=\{z\in\Z^2:\|z-x\|_\infty\le r\},\qquad
 B_r^\circ(x)=\{z\in\Z^2:\|z-x\|_\infty<r\},\\
 \partial B_r(x)=B_r(x)\setminus B_r^\circ(x).
 \end{gathered}
\]
Conditional integration over a square means integration over its open square
\(B_r^\circ(x)\).  In the pure model its exterior enters through
\(\partial B_r(x)\).  Enlarging an edit square by a fixed number of lattice
steps covers any alternate convention for this outer layer.

\subsection{The finite-state mechanism}

The conditional expectation in the scale flow has an exact probabilistic
meaning.  On a finite spin space, write
$d\mu_t=Z_t^{-1}e^{H_t}\,d\mu_{\rm ref}$, with a fixed strictly positive
reference law.  Fix update sets $I_x$ and bounded real functions $f_x$,
and let $P_{t,x}$ condition under $\mu_t$ on the complement of $I_x$.
If the log-density path satisfies
\[
 \dot H_t=c\sum_x(P_{t,x}f_x-f_x),\qquad c\ge0,
\]
then each summand has zero $\mu_t$-mean.  Differentiating $Z_t$ therefore
gives the probability derivative in the following lemma.  Adding a
constant to $\dot H_t$ gives the same normalized derivative.  The lemma
shows why moving a local function toward its conditional expectation
can be reversed by refreshes at rates independent of the current spins.

\begin{lemma}[Reverse generator identity]\label{lem:flow-reverse}
Let \(\mu_t\) be a differentiable family of strictly positive probability
laws on a finite spin space, and let \(I_x\) be specified update sets.
Let \(f_x\) be fixed bounded real functions, with
\(\|f_x\|_\infty\le a\), \(a\ge0\), and let \(P_{t,x}\) denote conditional
expectation under \(\mu_t\) given the complement of \(I_x\).  Suppose
\begin{equation}\label{eq:flow-exact-finite-derivative}
 \dot\mu_t(\sigma)=\mu_t(\sigma)c
 \sum_x\bigl(P_{t,x}f_x(\sigma)-f_x(\sigma)\bigr),
 \qquad c\ge0.
\end{equation}
If \(a>0\), define the refresh law
\begin{equation}\label{eq:flow-refresh-law}
 Q_{t,x}(d\eta_{I_x}\mid\sigma_{I_x^c})
 =\mu_t(d\eta_{I_x}\mid\sigma_{I_x^c})
 \left[1+\frac{f_x(\eta_{I_x}\sigma_{I_x^c})
                  -P_{t,x}f_x(\sigma)}{2a}\right].
\end{equation}
It is a probability law.  Updating set \(I_x\) from this law at rate
\(2ac\), independently for each \(x\), has generator \(\mathcal L_t\)
with
\(\mathcal L_t^*\mu_t=-\dot\mu_t\).  Thus using \(\mathcal L_{1-t}\)
from time zero to one carries the law at time one exactly to the law at
time zero.
If \(a=0\), the derivative in
\eqref{eq:flow-exact-finite-derivative} vanishes and no refresh is needed.
\end{lemma}

\begin{proof}
Since both \(f_x\) and its conditional expectation lie in \([-a,a]\),
the bracket in \eqref{eq:flow-refresh-law} lies in \([0,2]\).
Its conditional integral equals one.  For a configuration \(\sigma\),
the incoming mass of an \(x\)-refresh applied to \(\mu_t\) is
\[
 \mu_t(\sigma)\left[1+
 \frac{f_x(\sigma)-P_{t,x}f_x(\sigma)}{2a}\right].
\]
Indeed, summing the old interior configuration first leaves its exterior
marginal, which is then multiplied by the refresh conditional law.
Subtract the outgoing mass \(\mu_t(\sigma)\) and multiply by \(2ac\).
The result is
\(c\mu_t(\sigma)(f_x-P_{t,x}f_x)\).  Summing over \(x\) proves the
adjoint identity, including its sign and normalization.  The finite-state
forward equation has a unique solution, so the reversed family is the
law of the stated time-inhomogeneous chain.  The zero case is immediate.
\end{proof}

We will realize this operation on translation-invariant bulk interaction
densities at lattice scale $s$, with $c=s^{-2}$ and update squares of
radius $Ls$ for a fixed scale ratio $L$.
The next construction controls the new interactions generated by their
conditional expectations.  Its periodic evaluation has a small residual
relative to the finite-state identity; that residual will be estimated
before the lemma is used on tori.

\subsection{Spaces of labeled potentials}

An interaction created by conditional integration need not remain supported
in one square.  We retain a covering list as part of each coefficient, rather
than estimating an interaction only by the diameter of its support.

For an integer \(s\ge1\), a rooted list of length \(m\ge1\) is an ordered
tuple
\[
 \ell=(z_1,\ldots,z_m),\qquad z_1=0,\quad z_i\in\Z^2,
\]
such that the graph joining \(i\) and \(j\) when
\(B_s(z_i)\cap B_s(z_j)\ne\varnothing\) is connected.  Repetitions of
centers are allowed.  Put \(S_s(\ell)=\bigcup_iB_s(z_i)\).  A coefficient
on \(\ell\) is a function of the spins in \(S_s(\ell)\), modulo constants.
Whenever a representative is needed, we use the bounded linear choice
\[
 [v]\longmapsto v-\E_0v.
\]
If \(\qnorm{v}=\inf_{c\in\C}\|v-c\|_\infty\), this representative has
norm at most \(2\qnorm{v}\).  Thus working with these representatives gives
an equivalent norm to the quotient norm used in
Theorem~\ref{thm:transfer}.  The same choice works for complex coefficients.

For \(b>0\), let \(\mathcal A_{s,b}\) be the space of families
\(V=(v_\ell)_\ell\) of these centered coefficients with
\begin{equation}\label{eq:flow-list-norm}
 \|V\|_{s,b}:=\sum_\ell e^{b(|\ell|-1)}\|v_\ell\|_\infty<\infty.
\end{equation}
It is a Banach space: it is a weighted \(\ell^1\) sum of finite-dimensional
spaces.  For each fixed list length there are only finitely many lists,
since every center is within \(2s(m-1)\) of the root.  We use its closed
subspace invariant under spin flip and the square group; the group acts on
the lists and their functions together.  In particular its length-one
coefficient, denoted \(f=\pi_sV\), is an even square-invariant function on
\(B_s\).  Different lists may describe the same function.  No independence
of these descriptions is asserted or required.

The translation convention is fixed by the root: if
\((\tau_x\sigma)_z=\sigma_{x+z}\), the log-density addition represented by
\(V\) is
\begin{equation}\label{eq:flow-density}
 \mathcal H_s(V)=s^{-2}\sum_{x\in\Z^2}\sum_\ell
 v_\ell(\tau_x\sigma).
\end{equation}
Only conditional differences, or the corresponding finite torus sum, are
used; the full-plane sum itself is not assigned a numerical value.
Recentring a local coefficient by a lattice translation leaves this
translation sum unchanged.  We will not need such recentering: our changes
of covering lists always keep the root at zero.  Additive constants in the
coefficients likewise have no effect on any normalized measure.

Fix an integer \(L\ge9\).  For an exterior configuration \(\eta\) outside
\(I=B_{Ls}^\circ\), let
\begin{equation}\label{eq:flow-local-hamiltonian}
 H_{I,s}(V)(\sigma_I,\eta)
 =s^{-2}\sum_\ell\sum_{x:(x+S_s(\ell))\cap I\ne\varnothing}
 v_\ell(\tau_x(\sigma_I\eta)).
\end{equation}
Including terms that happen to be constant in \(\sigma_I\) is harmless.
For a fixed list of length \(m\), the number of translations in this sum
is at most \(C_Lms^2\): for at least one of its boxes to meet \(I\), its
center must lie in \(B_{(L+1)s}\), which has at most \(C_Ls^2\) sites.
Consequently
\begin{equation}\label{eq:flow-placement}
 \|H_{I,s}(V)\|_\infty
 \le C_L\sum_\ell |\ell|\|v_\ell\|_\infty.
\end{equation}
The constants in this and subsequent one-step estimates do not depend on
\(s\).

Let \(P^0_{Ls}\) be the pure conditional expectation over \(I\), and define
\begin{equation}\label{eq:flow-conditional}
 P^V_{Ls}f(\eta)=
 \frac{P^0_{Ls}\bigl(f\exp H_{I,s}(V)\bigr)(\eta)}
 {P^0_{Ls}\bigl(\exp H_{I,s}(V)\bigr)(\eta)}.
\end{equation}
For real \(V\), this is the conditional expectation in the perturbed
specification.  For complex \(V\) with a sufficiently small right-hand
side in \eqref{eq:flow-placement}, the denominator is uniformly separated
from zero.  Indeed its distance from \(1\) is at most
\(\exp(\|H_{I,s}(V)\|_\infty)-1\).  Equation
\eqref{eq:flow-conditional} is then defined for every exterior spin
configuration, including configurations of zero full-plane probability.

\subsection{The conditional differential equation}

During one step we hold the initial single-square function \(f\) fixed and
solve, modulo constants,
\begin{equation}\label{eq:flow-ode}
 \frac{d}{dt}V(t)=-f+P^{V(t)}_{Ls}f,
 \qquad V(0)=V,\qquad 0\le t\le1.
\end{equation}
The right-hand side is interpreted as a labeled family by the expansion
below.  This is the translation-density version of the finite-state
logarithmic derivative above: applying \(\mathcal H_s\) sums its translates
with coefficient \(s^{-2}\).  Holding \(f=\pi_sV\) fixed is essential:
it cancels the original single-square coefficient exactly at time one.

\begin{lemma}[Analytic conditional flow]\label{lem:flow-analytic}
Fix \(L\ge9\) and \(d>0\).  There are \(r>0\) and \(C<\infty\),
independent of the integer \(s\), such that for
\(\|V\|_{s,d}<r\), Equation~\eqref{eq:flow-ode} has a specified analytic
solution taking values in \(\mathcal A_{s,d/2}\).  Uniformly in
\(0\le t\le1\),
\begin{equation}\label{eq:flow-intermediate}
 \sum_\ell |\ell|e^{(d/2)(|\ell|-1)}
 \|v_\ell(t)\|_\infty\le C\|V\|_{s,d}.
\end{equation}
Its endpoint has the form
\begin{equation}\label{eq:flow-before-compression}
 V(1)=(V-f)+P^0_{Ls}f+\mathcal R_s(V),\qquad
 \|\mathcal R_s(V)\|_{s,d/2}\le C\|V\|_{s,d}^2.
\end{equation}
On a smaller ball the remainder is analytic and satisfies
\(\|D\mathcal R_s(V)\|\le C\|V\|_{s,d}\) and
\(\|D^2\mathcal R_s(V)\|\le C\), with these operator norms from
\(\mathcal A_{s,d}\) to \(\mathcal A_{s,d/2}\).
The construction preserves real coefficients and the stated symmetries.
\end{lemma}

\begin{proof}
Choose a connected covering list \(\mathcal C\) of \(B_{Ls}\) by
\(c_L\ge2\) squares of radius \(s\), all contained in \(B_{Ls}\), including
the root square.  Such a list can be obtained from centers
\((as,bs)\), \(-L+1\le a,b\le L-1\).  It has size depending only on
\(L\), and its union is exactly \(B_{Ls}\).

Expand the numerator in \eqref{eq:flow-conditional} in its exponential
series.  Expand the reciprocal denominator as
\[
 \frac1{P^0e^H}
 =\sum_{j\ge0}\bigl(-P^0(e^H-1)\bigr)^j.
\]
The absolute scalar series is bounded by
\[
 \|f\|_\infty\frac{e^h}{2-e^h},\qquad
 h=\|H\|_\infty<\log2.
\]
Its coefficient of degree \(p\) is bounded by \(K_0K_1^p\|f\|_\infty\)
for fixed numerical \(K_0,K_1\).  Each degree-\(p\) coefficient is a sum
of products of pure conditional expectations with precisely \(p\)
inserted translates of coefficients of \(V\).  It depends only on the
spins in the covering list \(\mathcal C\) and the inserted lists.  Give it
the concatenated list consisting of \(\mathcal C\), followed by all
inserted lists, retaining repeated boxes.  An inserted list meets the
integrated square and hence is connected to \(\mathcal C\).  If its input
lengths are \(m_1,\ldots,m_p\), its output length is exactly
\begin{equation}\label{eq:flow-label-size}
 n=c_L+\sum_{i=1}^p m_i.
\end{equation}
Center each output coefficient with the fixed linear projection.  This
costs a factor of at most two in its sup norm.  The placement bound
\eqref{eq:flow-placement} supplies a factor \(C_Lm_i\) per insertion;
the factors \(s^2\) have canceled with the translation normalization.
Taking the square-group average of this rule makes it equivariant, without
increasing a sum of absolute norms.  Spin flip commutes with every pure
conditional integration.

Here is a convergent majorant, including the dependence on list size.
Let \(a_m(t)\) be the sum of nonnegative coefficient majorants on lists of
length \(m\), and put
\[
 X_b(t)=\sum_{m\ge1}m e^{b(m-1)}a_m(t),\qquad
 Y_b(t)=\sum_{m\ge1}m(m-1)e^{b(m-1)}a_m(t).
\]
The majorant includes the absolute values of the initial coefficients and
of both fixed forcing terms \(-f\) and \(P^0f\).  By
\eqref{eq:flow-label-size}, for \(d/2\le b\le d\) the contribution of
terms with \(p\ge1\) insertions to its \(X_b\) derivative is at most
\begin{equation}\label{eq:flow-majorant-term}
 K\|f\|_\infty K^p
 \left[c_L X_b^p+p(X_b+Y_b)X_b^{p-1}\right],
\end{equation}
where \(K=K_{L,d}\).  To verify this bound, the exponential weight of an
output is the product of the input weights times
\(e^{b(c_L-1+p)}\), absorbed in \(K^{p+1}\).  Its additional size factor
\(n\) is the sum of \(c_L\) and the \(m_i\).  Choosing the latter term
at the \(i\)-th insertion replaces one factor
\(\sum m e^{b(m-1)}a_m=X_b\) by
\(\sum m^2 e^{b(m-1)}a_m=X_b+Y_b\).  This proves
\eqref{eq:flow-majorant-term}, including the unbounded size factor.

Set \(b(t)=3d/4-dt/8\), reserving a final exponent gap
\(b(1)-d/2=d/8\).  As long as \(KX_{b(t)}\le1/2\), summing the
geometric series and its derivative in \eqref{eq:flow-majorant-term}
gives
\begin{equation}\label{eq:flow-majorant-differential}
 D^+X_{b(t)}(t)
 \le -\frac d8Y_{b(t)}(t)
       +C_{L,d}\|f\|_\infty\bigl(1+Y_{b(t)}(t)\bigr).
\end{equation}
Furthermore
\[
 X_{3d/4}(0)\le
 \left(\sup_{m\ge1}m e^{-d(m-1)/4}\right)\|V\|_{s,d}.
\]
Choose the input radius so that \(C_{L,d}\|f\|_\infty\le d/16\) and
the resulting bound
\(X_{b(t)}(t)\le X_{3d/4}(0)+C_{L,d}\|f\|_\infty t\) is strictly
below \(1/(2K)\).  A first-exit argument then gives this bound through
time one.  In particular it gives \eqref{eq:flow-intermediate}.

For completeness, the same estimates construct the solution rather than
only bounding a hypothetical one.  First retain only output lists of
length at most \(M\).  By \eqref{eq:flow-label-size}, every nonlinear
coefficient equation at length \(n\) uses only coefficients of lengths
strictly smaller than \(n\); only the fixed forcing occurs at length one
or at the fixed cover length.  The truncated equations therefore have
successive explicit integral solutions.  Their coefficients are
polynomials in the finitely many input coefficients of lengths at most
\(M\).  Truncations agree on all smaller lengths.  The corresponding
nonnegative triangular majorants obey
\eqref{eq:flow-majorant-differential} and increase to a majorant for all
coefficients.  Their tails satisfy, uniformly on the complex input ball,
\[
 \sup_{0\le t\le1}\sum_{m>M}e^{(d/2)(m-1)}a_m(t)
 \le \frac{C\|V\|_{s,d}}{M+1}.
\]
Thus the truncations converge uniformly into
\(C([0,1],\mathcal A_{s,d/2})\).  They are holomorphic polynomials in
finite-dimensional projections of the input, so their locally uniform
limit is analytic.  The reserved exponent gap also gives uniform
convergence of the derivatives in \(\mathcal A_{s,d/2}\): the vector
field is bounded in \(\mathcal A_{s,5d/8}\) by
\(C\|f\|_\infty[1+\sum_{p\ge1}(KX_{5d/8})^p]\), so its tail above
length \(M\), measured at exponent \(d/2\), is at most
\(C\|f\|_\infty e^{-dM/8}\).  Absolute convergence permits termwise
conditional integration and time differentiation.  This proves
\eqref{eq:flow-ode}; uniqueness for this labeled expansion follows
successively in each length.  This construction is also valid for real
data and commutes with symmetry averaging.

Integrating the two fixed forcing terms gives
\((V-f)+P^0f\) at time one.  Every remaining term contains one factor
\(f\) and at least one coefficient of \(V(t)\).  Its
\(\mathcal A_{s,d/2}\) norm is bounded by
\(C\|f\|_\infty\sum_{p\ge1}(KX_{d/2}(t))^p\), which is
\(O(\|V\|_{s,d}^2)\).  Integration proves
\eqref{eq:flow-before-compression}.  The same analytic bound on a complex
ball, followed by the Cauchy estimates on a smaller ball, gives the two
derivative estimates.  All constants used above are independent of \(s\).
\end{proof}

The loss of exponential weight has a specific purpose: it pays for the
number of placements of a long interaction.  The next geometric step
restores that weight when the scale is increased.

\begin{lemma}[Compression of covering lists]\label{lem:flow-compression}
For every integer \(L\ge9\), there is a linear relabeling map
\(C_s:\mathcal A_{s,d/2}\to\mathcal A_{Ls,d}\) that leaves the local
functions and their roots unchanged and satisfies
\[
 \|C_sW\|_{Ls,d}\le\|W\|_{s,d/2}.
\]
If a declared union is contained in \(B_{Ls}\), it is relabeled as a
single square.  The map commutes with the specified symmetries.  The
rewriting map \(\mathcal S_s=L^2C_s\) preserves the translation density:
\begin{equation}\label{eq:flow-rewrite}
 \mathcal H_{Ls}(\mathcal S_sW)=\mathcal H_s(W).
\end{equation}
On input lists of length at least two it also satisfies
\begin{equation}\label{eq:flow-carry}
 \|\mathcal S_sW\|_{Ls,d}
 \le L^2 e^{-d/2}\|W\|_{s,d}.
\end{equation}
\end{lemma}

\begin{proof}
Choose a spanning tree of the connected box graph using its vertex labels,
and traverse its doubled edges starting at the root.  The resulting walk
has \(2(m-1)\) steps, each of sup-norm length at most \(2s\).  Put
\(q=\lfloor(L-1)/2\rfloor\), and retain the centers at times
\(0,q,2q,\ldots\), through the last such time in the walk.  Every old
center occurs in the walk and is at distance at most \(2s(q-1)\) from
the preceding retained center.  Its radius-\(s\) square is therefore
contained in the radius-\(Ls\) square at that retained center.  Consecutive
retained squares overlap.  Removing repetitions leaves a connected
covering list, containing the root, of length \(M\) with
\[
 M-1\le\left\lfloor\frac{2(m-1)}q\right\rfloor
 \le\frac{m-1}{2}.
\]
In the final inequality we used \(q\ge4\).  If the original union fits
in \(B_{Ls}\), choose that single square instead.  The same inequalities
hold.  The procedure uses only the labeled graph, so it can be chosen
square-equivariantly; alternatively one can average its conjugates under
the square group.  Such averaging preserves all norm bounds.

Relabeling a function does not change it.  The inequality
\(e^{d(M-1)}\le e^{(d/2)(m-1)}\), summed over lists, proves the first
bound.  For \(m\ge2\), comparison with the original weight gives the
extra factor \(e^{-d(m-1)/2}\le e^{-d/2}\), proving
\eqref{eq:flow-carry}.  Finally, the factor \(L^2\) cancels the change
from \(s^{-2}\) to \((Ls)^{-2}\) in the translation normalization, which
proves \eqref{eq:flow-rewrite} term by term.
\end{proof}

Define the one-step map \(\Phi_s(V)=\mathcal S_sV(1)\).  Its linearization
is
\begin{equation}\label{eq:flow-linearization}
 T_sV=A_sf+\mathcal S_s(V-f),\qquad
 A_sf=L^2P^0_{Ls}f,\qquad f=\pi_sV,
\end{equation}
where the first term is a single-square coefficient at scale \(Ls\).
Indeed the pure conditional term is supported in the fixed cover of
\(B_{Ls}\), so the special single-square rule of
Lemma~\ref{lem:flow-compression} applies.  Lemmas~\ref{lem:flow-analytic}
and~\ref{lem:flow-compression} give
\begin{equation}\label{eq:flow-map-remainder}
 \Phi_s(V)=T_sV+R_s(V),\qquad
 \|R_s(V)\|_{Ls,d}\le C\|V\|_{s,d}^2,
 \qquad \|DR_s(V)\|\le C\|V\|_{s,d},
\end{equation}
uniformly in \(s\), on a common ball.  The large, but fixed, factor
\(L^2\) is included in \(C\).

\subsection{A decaying graph for the nonautonomous maps}

We next solve the recurrence \(V_{k+1}=\Phi_{s_k}(V_k)\).  The spaces and
maps vary with \(k\), so we give the sequence-space argument rather than
appeal to an autonomous stable-manifold theorem.

Write \(s_k=L^ks_0\) and \(\mathcal X_k=\mathcal A_{s_k,d}\).
Extend the functional \(\ell_{s_k}\) of
Theorem~\ref{thm:transfer} to \(\mathcal X_k\) by
\(\ell_kV=\ell_{s_k}(\pi_{s_k}V)\).  The same theorem, and the equivalence
of our norm with the quotient norm, allow the following choices.  First
choose its target error \(a>0\) sufficiently small, then \(L\) sufficiently
large, then the farther measurement ratio \(H\), approximation accuracies,
and \(s_0\) in the order prescribed there.  All sufficiently large
\(s\), including every \(s_k\), then satisfy
\begin{equation}\label{eq:flow-transfer-use}
 \|A_sf\|_\infty\le C_1L|\ell_sf|+a\|f\|_\infty,
 \qquad
 |\ell_{Ls}A_sf-L\ell_sf|\le a\|f\|_\infty.
\end{equation}
Constants have been enlarged, and the input error decreased, to account
for centering representatives.  The norms of \(\ell_s\) and the norms
of some right inverses are bounded independently of \(s\) and \(L\).
After fixing \(L\), take \(d\) large enough that
\(\kappa=L^2e^{-d/2}\) is as small as needed.  Increasing \(s_0\) further
does not affect any of these bounds.  This freedom will permit the fixed
pure annulus estimates in Section~\ref{sec:stability} to be imposed before
we construct the trajectory.

\begin{lemma}[Decaying sequence graph]\label{lem:flow-stable-graph}
With the choices just described, there are bounded single-square vectors
\(w_k\in\mathcal X_k\), \(\ell_kw_k=1\), and a real analytic function
\(h\) on a neighborhood of zero in \(\ker\ell_0\), with \(h(0)=0\),
such that each sufficiently small vector
\begin{equation}\label{eq:flow-graph}
 V_0=h(u_0)w_0+u_0,\qquad u_0\in\ker\ell_0,
\end{equation}
has a trajectory \(V_{k+1}=\Phi_{s_k}(V_k)\) satisfying
\[
 \sup_{k\ge0}2^k\|V_k\|_{s_k,d}\le C\|u_0\|_{s_0,d}.
\]
In a sufficiently small ball these are exactly the initial conditions
with a trajectory of finite sufficiently small norm
\(\sup_k2^k\|V_k\|_{s_k,d}\).  The tangent space of the graph at zero is
exactly the set of initial conditions whose linear trajectories
\(Z_{k+1}=T_{s_k}Z_k\) have finite norm
\(\sup_k2^k\|Z_k\|_{s_k,d}\).
\end{lemma}

\begin{proof}
Take a bounded single-square right inverse \(w_0\).  Recursively define
\[
 a_k=\ell_{k+1}T_{s_k}w_k,\qquad
 w_{k+1}=a_k^{-1}T_{s_k}w_k.
\]
The vectors remain single-square by \eqref{eq:flow-linearization}.  If
\(\|w_k\|\le M\), Equation~\eqref{eq:flow-transfer-use} gives
\[
 |a_k-L|\le aM,\qquad
 \|w_{k+1}\|\le\frac{C_1L+aM}{L-aM}.
\]
Choose a fixed \(M\) larger than the initial right-inverse bound and
\(2C_1+1\), and choose the parameters to make the displayed ratio at most
\(M\).  Induction gives \(\sup_k\|w_k\|\le M\) and
\(a_k\ge L/2\ge4\).  We have used real vectors, so the scalar multipliers
are positive; their complexifications use the same fixed multipliers.

Every \(V_k\) has the unique decomposition
\(V_k=\alpha_kw_k+u_k\), with \(\alpha_k=\ell_kV_k\) and
\(u_k\in\ker\ell_k\).  These coordinates have norms uniformly
equivalent to \(\|V_k\|\).  For \(u\in\ker\ell_k\), its single-square
part has \(\ell_{s_k}\pi_{s_k}u=0\).  Consequently
\eqref{eq:flow-transfer-use} and \eqref{eq:flow-carry} imply
\[
 \|T_{s_k}u\|\le(a+\kappa)\|u\|.
\]
Set
\[
 b_ku=\ell_{k+1}T_{s_k}u,\qquad
 D_ku=T_{s_k}u-w_{k+1}b_ku.
\]
Thus \(D_k:\ker\ell_k\to\ker\ell_{k+1}\).  The uniform functional
and vector bounds give \(\|b_k\|+\|D_k\|\le C(a+\kappa)\).
Choose \(a\) and then \(\kappa\) so that \(\|D_k\|\le1/8\).
The coordinate form of the linear recurrence is exactly triangular:
\begin{equation}\label{eq:flow-triangular}
 \alpha_{k+1}=a_k\alpha_k+b_ku_k,
 \qquad u_{k+1}=D_ku_k.
\end{equation}
There is no term from \(\alpha_k\) into the second equation, by the
definition of \(w_{k+1}\).

We record the inverse of the inhomogeneous recurrence.  Given \(u_0\)
and forcing terms \(g_k\in\C\), \(v_k\in\ker\ell_{k+1}\) with finite
norm \(\sup_k2^k\max\{|g_k|,\|v_k\|\}\), its only
decaying solution is
\begin{align}
 u_k&=D_{k-1}\cdots D_0u_0
       +\sum_{j=0}^{k-1}D_{k-1}\cdots D_{j+1}v_j,
       \label{eq:flow-forward-sum}\\
 \alpha_k&=-\sum_{j=k}^{\infty}
       \frac{b_ju_j+g_j}{a_k a_{k+1}\cdots a_j}.
       \label{eq:flow-backward-sum}
\end{align}
An empty product in these expressions is the identity.  The first formula
is ordinary forward substitution.  Iterating the scalar equation to time
\(n\), dividing by \(a_k\cdots a_n\), and letting \(n\to\infty\)
gives the second; the terminal term vanishes for a decaying sequence.
Conversely, absolutely convergent sums in
\eqref{eq:flow-forward-sum}--\eqref{eq:flow-backward-sum} satisfy both
equations.  In the norm
\[
 \|(\alpha,u)\|_*=
 \sup_{k\ge0}2^k\max\{|\alpha_k|,\|u_k\|\},
\]
they are a bounded linear function of \(u_0\) and the forcing norm
\(\sup_k2^k\max\{|g_k|,\|v_k\|\}\).  For example the forward
convolution is bounded by the geometric sum
\(2\sum_{j\ge0}(2/8)^j\), and the backward convolution by
\(\sum_{j\ge0}4^{-j-1}2^{-j}\).  These constants are uniform in \(k\)
and in the lattice scales.

For the nonlinear recurrence use the forcing
\[
 g_k=\ell_{k+1}R_{s_k}(\alpha_kw_k+u_k),\qquad
 v_k=(I-w_{k+1}\ell_{k+1})R_{s_k}(\alpha_kw_k+u_k).
\]
Equation~\eqref{eq:flow-map-remainder} implies that on
\(\|(\alpha,u)\|_*\le\rho\) this forcing has norm at most \(C\rho^2\)
and Lipschitz constant at most \(C\rho\).  To see the latter explicitly,
at time \(k\) both arguments have norm at most \(C\rho2^{-k}\), and
their difference has norm at most
\(C2^{-k}\|(\alpha,u)-(\widetilde\alpha,\widetilde u)\|_*\).
Multiplication by \(2^k\) leaves a factor at most \(C\rho\).

Insert this forcing into the two sums.  Their bounded linear inverse is a
contraction on a ball of radius \(C\|u_0\|\) when \(u_0\) is small
enough.  It produces the asserted unique sequence.  Iteration of the
contraction also proves analytic dependence on complex \(u_0\): the
iterates are analytic and converge uniformly on smaller balls.  Set
\(h(u_0)=\alpha_0\).  The construction respects conjugation, so \(h\)
is real analytic on real data.

Finally the derivative at zero is given by the same sums with the
quadratic forcing deleted.  It parametrizes exactly all decaying
solutions of \eqref{eq:flow-triangular}.  Uniform equivalence of
coordinates proves the assertion about the tangent space and the
trajectory estimate.
\end{proof}

\subsection{Energy transversality and the choice of temperature}

We apply the graph to the physical interaction.  To fix its normalization,
let
\[
 E(\sigma)=\frac12\sum_{z:\|z\|_1=1}\sigma_0\sigma_z.
\]
Its translation sum counts each nearest-neighbor bond once.  By
Lemma~\ref{lem:contour-spin}, the interaction from \(U\) has a bounded
finite-support even square-invariant density \(u_U\) such that its
translation sum is the full sum of disagreement monomials, up to
exterior-only constants when spins are pinned.  One direct choice is to
sum \(U(S)\prod_{e\in S}d_e(\sigma)\), with \(d_e\) as in
Section~\ref{sec:introduction}, over one
representative of each translation class of nonempty edge sets, then
average the density over the square group.  There are finitely many such
classes because \(U\) has finite range.  Translation invariance shows
that this averaging preserves the translation sum exactly.  Enlarge
\(s_0\), if needed, so that all these representatives are supported in
\(B_{s_0}\).

Relative to the pure interaction at
\(\beta_0=\log(1+\sqrt2)/(2J)\), the initial potential is the
single-square vector
\begin{equation}\label{eq:flow-physical-input}
 V_0(\beta,\lambda)=s_0^2
 \left[J(\beta-\beta_0)E-\beta\lambda u_U\right],
\end{equation}
with its additive constant removed.  The minus sign agrees with the
contour penalty.  In particular this formula retains every interaction
crossing a domain boundary: it is a bulk density whose conditional law
is the law in Lemma~\ref{lem:contour-spin}, not a density obtained by
discarding boundary terms.

\begin{theorem}[Domain-independent decaying trajectory]
\label{thm:critical-trajectory}
For every \(J>0\) and finite-range real potential \(U\) satisfying the
symmetries of Theorem~\ref{thm:main}, the preceding scale parameters can
be fixed independently of all domains and boundary data.  The initial
scale \(s_0\) may be enlarged by any prescribed finite amount before the
following construction.  For each sufficiently small \(\epsilon>0\)
there is \(\lambda_*>0\) and a real analytic function
\(\beta_c:(-\lambda_*,\lambda_*)\to(0,\infty)\), with
\(\beta_c(0)=\beta_0\), for which the initial potential
\eqref{eq:flow-physical-input} has a trajectory
\(V_{k+1}=\Phi_{s_k}(V_k)\) satisfying
\begin{equation}\label{eq:flow-critical-decay}
 \|V_k\|_{s_k,d}\le\epsilon2^{-k},\qquad k\ge0.
\end{equation}
Its intermediate one-step trajectories obey
\begin{equation}\label{eq:flow-critical-intermediate}
 \sup_{0\le t\le1}
 \sum_\ell |\ell|e^{(d/2)(|\ell|-1)}
 \|v_{k,\ell}(t)\|_\infty\le C\epsilon2^{-k}.
\end{equation}
The same construction applies to both signs of \(\lambda\).
\end{theorem}

\begin{proof}
Let \(G(V)=\ell_0V-h((I-w_0\ell_0)V)\).  The graph in
Lemma~\ref{lem:flow-stable-graph} is exactly \(G(V)=0\).  We prove that
the derivative of \(G(V_0(\beta,\lambda))\) in \(\beta\) at
\((\beta_0,0)\) is nonzero.

The temperature derivative of \eqref{eq:flow-physical-input} is the
single-square vector \(Z_0=Js_0^2E\), modulo constants.  Its successive
linear images remain single-square and are
\begin{equation}\label{eq:flow-energy-iterate}
 Z_k=Js_k^2\E_0[E\mid\sigma_{\partial B_{s_k}}],\qquad k\ge1,
\end{equation}
modulo constants.  For the induction step, the pure Markov property gives
\[
 \E_0\!\left[\E_0[E\mid\sigma_{\partial B_{s_k}}]
              \mid\sigma_{\partial B_{s_{k+1}}}\right]
 =\E_0[E\mid\sigma_{\partial B_{s_{k+1}}}].
\]
Although the two boundary sigma-fields are not nested, this identity
follows by first conditioning on the full exterior of \(B_{s_k}^\circ\).
The larger boundary is measurable there, and the conditional expectation
of \(E\) given that exterior is the inner-boundary conditional expectation.
Multiplication by \(L^2\) at each step proves
\eqref{eq:flow-energy-iterate}.

Put \(y_k=(2s_k,0)\) and \(E_{y_k}=E\circ\tau_{y_k}\).  This bounded
observable is exterior to \(B_{s_k}^\circ\), so another use of the pure
Markov property yields
\begin{equation}\label{eq:flow-energy-test}
 \Cov_0(Z_k,E_{y_k})
 =Js_k^2\Cov_0(E,E_{y_k}).
\end{equation}
Proposition~\ref{prop:energy} says that the right-hand side converges to
a strictly positive number.  Its square symmetrization is precisely what
rules out cancellation between bond orientations.  On the other hand
\[
 |\Cov_0(Z_k,E_{y_k})|
 \le 2\|E\|_\infty\qnorm{Z_k}.
\]
Thus the quotient norms of \(Z_k\) do not tend to zero.  If \(Z_0\)
were tangent to \(G=0\), Lemma~\ref{lem:flow-stable-graph} would imply
\(\|Z_k\|\le C2^{-k}\), a contradiction.  This proves the required
nonzero derivative.

The scalar analytic implicit-function theorem applied to
\(G(V_0(\beta,\lambda))=0\) now gives a unique local real analytic branch
\(\beta=\beta_c(\lambda)\) through \((\beta_0,0)\).  All data are real,
so the branch is real on a two-sided real interval.  Shrink the interval
so that \(\beta_c(\lambda)>\beta_0/2>0\) and the initial point lies in
the small graph neighborhood.  The analytic sequence constructed in
Lemma~\ref{lem:flow-stable-graph} vanishes at \(\lambda=0\), hence its
sequence norm is at most \(C_U|\lambda|\) on a smaller interval.  Shrink
it once more to obtain \eqref{eq:flow-critical-decay}.  Equation
\eqref{eq:flow-critical-intermediate} is then
Lemma~\ref{lem:flow-analytic}.  Every ingredient in this construction is
a bulk object fixed before a domain or exterior configuration is chosen.
\end{proof}

The notation \(\beta_c\) refers to this constructed temperature.  The
remaining sections establish its claimed curve-scaling property.  No
sign condition on \(u_U\), or on \(\lambda\), was used to construct it.

\subsection{Reverse refreshes and independent marked squares}

For an integer \(N\), let \(\mathbb T_N=(\Z/N\Z)^2\), and write
\(\mu_{0,N}\) for its pure critical spin law.  Evaluating local functions
on periodic spin configurations defines the finite Hamiltonian
\[
 \mathcal H_{s,N}(V)(\sigma)
 =s^{-2}\sum_{x\in\mathbb T_N}\sum_\ell v_\ell(\tau_x\sigma),
 \qquad
 d\mu_{s,N}^{V}=\frac{\exp\mathcal H_{s,N}(V)}
 {\E_{0,N}\exp\mathcal H_{s,N}(V)}\,d\mu_{0,N}.
\]
These sums converge absolutely in sup norm for the potentials above.
For \(V_0\) from Theorem~\ref{thm:critical-trajectory}, the law
\(\mu_{s_0,N}^{V_0}\) is the physical perturbed torus law, denoted
\(\mu_{\lambda,N}\), for all sufficiently large \(N\).

The scale index $k$ labels successive interaction densities; within each
step, $t\in[0,1]$ is the time variable.  We will use the refresh laws of
Lemma~\ref{lem:flow-reverse} on each path, with the following estimate
controlling the error in its probability equation.

Our flow was constructed in the plane, so its periodic evaluation need
not solve that finite-volume equation exactly.  We quantify the discrepancy
before applying the identity.

\begin{lemma}[Periodization error]\label{lem:flow-periodization}
Suppose \(Ls\le N/10\), and suppose a real one-step trajectory obeys
\[
 \sup_{0\le t\le1}\sum_\ell |\ell|
 e^{(d/2)(|\ell|-1)}\|v_\ell(t)\|_\infty\le r_s,
 \qquad \|f\|_\infty\le r_s,
\]
with \(r_s\) in the fixed small ball.  Let \(\mu_t=\mu_{s,N}^{V(t)}\),
and let \(P_{t,x}^{N}\) be its exact torus conditional expectation over
\(B_{Ls}^\circ(x)\).  There is a signed measure \(e_t\), of total mass
zero, such that
\begin{equation}\label{eq:flow-finite-residual}
 \dot\mu_t=\mu_t s^{-2}\sum_{x\in\mathbb T_N}
       (P_{t,x}^{N}f_x-f_x)+e_t,
 \qquad
 \|e_t\|_{\TV}\le C r_s(N/s)^2e^{-cN/s}.
\end{equation}
The constants \(C,c>0\) depend on the fixed flow parameters, but not on
\(s,N,t\).
\end{lemma}

\begin{proof}
A connected list of \(m\) radius-\(s\) squares has sup-norm diameter at
most \(2sm\).  Set \(M=\lfloor N/(8s)\rfloor\), changing fixed
constants if needed.  Every list of length less than \(M\) then has
diameter less than \(N/2\).  Such a translated support cannot intersect
two distinct periodic copies of \(B_{Ls}^\circ\): the distance between
these copies is at least \(N-2Ls\ge4N/5\), greater than that diameter.
Accordingly, for the potential truncated to these lists, the plane
conditional Hamiltonian with periodic exterior data is the same as the
torus conditional Hamiltonian, up to terms constant in the updated spins.
This is an identity of local Hamiltonians, not an assertion that plane and
torus Gibbs measures are equal.

For the discarded lists, the weighted estimate gives
\begin{equation}\label{eq:flow-tail-bound}
 \sum_{|\ell|\ge M}|\ell|\|v_\ell(t)\|_\infty
 \le C r_s e^{-cN/s}.
\end{equation}
The placement count of \eqref{eq:flow-placement} therefore bounds the
sup norm of their plane conditional Hamiltonian by
\(Cr_s e^{-cN/s}\).  The same bound holds for the torus conditional
Hamiltonian.  To see the latter, enumerate every torus translation that
can depend on the update spins by choosing one of its declared boxes
meeting a periodic copy of the update square.  The number of distinct
translations is at most \(C_Lms^2\); counting a translation more than
once only increases the upper bound.  This proves the torus version of
the same placement estimate, including supports with repeated periodic
images.

If two real conditional Hamiltonians differ by at most \(h\) after
subtracting constants, their normalized conditional densities have ratio
between \(e^{-2h}\) and \(e^{2h}\).  Thus their expectations of a bounded
function differ by at most \(C\|f\|_\infty h\) for \(h\le1\).
Apply this observation to truncation of each Hamiltonian.  The exact
identity for short lists and \eqref{eq:flow-tail-bound} yield, uniformly
in the entire spin configuration,
\begin{equation}\label{eq:flow-conditional-periodic-difference}
 \left|\left.(P_{Ls}^{V(t)}f)_x\right|_{\mathrm{periodic}}
             -P_{t,x}^{N}f_x\right|
 \le C r_s e^{-cN/s}.
\end{equation}
We used \(\|f\|_\infty\le r_s\le C\) to weaken the bound and simplify
notation.  For a long list, periodic evaluation of the plane conditional
function can involve an image of an updated spin in what was the plane
exterior.  This causes no problem: the comparison is pointwise, and the
entire dependence in question is included in the discarded tail.

Differentiate the finite Gibbs density and use
Equation~\eqref{eq:flow-ode} with periodic evaluation.  Replacing each
plane conditional term by its exact torus counterpart changes the
logarithmic derivative by at most
\(Cr_s(N/s)^2e^{-cN/s}\), before centering.  The torus counterpart has
zero \(\mu_t\)-mean.  Subtracting the mean of the error costs at most
another factor two.  Multiplying by \(\mu_t\) proves
\eqref{eq:flow-finite-residual}.
\end{proof}

\begin{proposition}[Coupling by independent marked squares]
\label{prop:edit-coupling}
Let the physical torus law \(\mu_{\lambda,N}\) be at the temperature of
Theorem~\ref{thm:critical-trajectory}, with trajectory bound
\eqref{eq:flow-critical-decay}.  Let
\[
 K_N=\max\{k\ge0:Ls_k\le N/10\}
\]
for all sufficiently large \(N\).  There is a coupling of
\(\sigma^{(0)}\sim\mu_{0,N}\),
\(\sigma^{(\lambda)}\sim\mu_{\lambda,N}\), and Bernoulli marks
\(M_{k,x}\), \(0\le k\le K_N\), \(x\in\mathbb T_N\), with the
following properties.  The marks are mutually independent and independent
of \(\sigma^{(0)}\), and
\begin{equation}\label{eq:flow-mark-probabilities}
 \P(M_{k,x}=1)\le C\epsilon2^{-k}s_k^{-2}.
\end{equation}
Except on an event of probability \(o(1)\) as \(N\to\infty\),
\(\sigma^{(\lambda)}\) can be obtained from \(\sigma^{(0)}\) by changing
spins only inside the union of the marked squares \(B_{Ls_k}(x)\).
The values of the changed spins may depend on all previous refreshes.
The constant and the exceptional probability are independent of test
domains and exterior configurations.  Enlarging these squares by any
fixed number of lattice steps preserves the conclusion.
\end{proposition}

\begin{proof}
For the duration of this proof put \(K=K_N\).  At each scale retain the
deterministic intermediate potential \(V_k(t)\) supplied by
Lemma~\ref{lem:flow-analytic}, with \(f_k=\pi_{s_k}V_k\) fixed.
The endpoint Gibbs law is exactly the starting law for the next scale:
\[
 \mu_{s_k,N}^{V_k(1)}=\mu_{s_{k+1},N}^{V_{k+1}},
\]
by the density identity \eqref{eq:flow-rewrite}, also after periodic
evaluation.

First bound the final perturbation.  Maximality of \(K\) gives
\(N/s_{K+1}<10L\), whereas \(K\to\infty\) with \(N\).  Hence
\begin{equation}\label{eq:flow-terminal-density}
 \|\mathcal H_{s_{K+1},N}(V_{K+1})\|_\infty
 \le (N/s_{K+1})^2\|V_{K+1}\|_{s_{K+1},d}
 \le C_L\epsilon2^{-K-1}=o(1).
\end{equation}
The normalized density relative to \(\mu_{0,N}\) lies between the
exponentials of plus and minus twice this bound.  Its total variation
distance from \(\mu_{0,N}\) therefore tends to zero.  Couple a sample
from this endpoint law and a pure sample so that they agree outside an
event of that probability.

For each site and each scale, now independently generate a Poisson process
on \([0,1]\) of rate
\[
 r_{k,x}=2s_k^{-2}\|f_k\|_\infty.
\]
These rates are deterministic; they do not involve the current spins,
the exterior data of a refresh, or its time within the step.  If \(f_k=0\)
we take no clocks at that scale.  Generate all these processes independently
of the coupled endpoint and pure samples, and take \(M_{k,x}=1\) exactly
when its process has at least one point.  Then
\[
 \P(M_{k,x}=1)=1-e^{-r_{k,x}}
 \le2s_k^{-2}\|f_k\|_\infty
 \le2\epsilon2^{-k}s_k^{-2},
\]
and all asserted independence properties already hold.  Additional
randomness used to choose refreshed spins can likewise be generated in
advance; no independence assertion about those spin values is needed.

Run scales in the order \(K,K-1,\ldots,0\).  At scale \(k\), traverse
the deterministic measure path backwards in time.  At a clock of site
\(x\), refresh \(B_{Ls_k}^\circ(x)\) with the law in
\eqref{eq:flow-refresh-law}, using the exact torus conditional of
\(\mu_{s_k,N}^{V_k(t)}\) and the fixed function \((f_k)_x\).  The
normalization and nonnegativity proved in Lemma~\ref{lem:flow-reverse}
apply at every time, irrespective of whether the evolving sample has the
instantaneous Gibbs law.  There are finitely many Poisson clocks on a
finite torus, so this constructs a finite-state time-inhomogeneous Markov
chain.  It changes spins only in marked squares.

Lemma~\ref{lem:flow-periodization}, with
\(r_s=C\epsilon2^{-k}\) from
\eqref{eq:flow-critical-intermediate}, bounds the error of its desired
reverse probability equation by
\[
 C\epsilon2^{-k}(N/s_k)^2e^{-cN/s_k}.
\]
These errors add in total variation.  Explicitly, for a prescribed
probability path \(p_t\) and a Markov evolution with generator
\(\mathcal L_t\), the difference satisfies the inhomogeneous forward
equation with forcing \(\dot p_t-\mathcal L_t^*p_t\).  Variation of
constants expresses the difference as the initial difference propagated
by a Markov kernel, plus the integral of the propagated forcing.
Markov kernels contract total variation of signed measures, so its norm
is bounded by the initial norm plus the integrated forcing norm.  Apply
this identity on each unit-time reverse step and use the exact equality
of consecutive endpoint laws.  The output law differs from
\(\mu_{\lambda,N}\) by at most
\begin{equation}\label{eq:flow-error-sum}
 C\sum_{k=0}^{K_N}\epsilon2^{-k}(N/s_k)^2e^{-cN/s_k}.
\end{equation}
This quantity tends to zero.  For each fixed \(k\), \(N/s_k\to\infty\).
Moreover \(r^2e^{-cr}\) is bounded for \(r\ge0\), so the summands,
extended by zero for \(k>K_N\), are dominated by a constant times the
summable sequence \(\epsilon2^{-k}\).  Dominated convergence proves
the assertion without requiring the largest-scale error to vanish at a
fixed radius ratio.

Finally maximally couple the output with a sample from the exact desired
law.  Its probability of disagreement is bounded by
\eqref{eq:flow-error-sum}.  This final correction can be sampled
conditionally on the already constructed output and leaves the joint law
of the pure input and all marks unchanged.  On the complement of the
endpoint-coupling failure and this final disagreement, the target sample
is obtained from the pure sample by precisely the described marked-square
refreshes.  Both exceptional probabilities tend to zero, which proves
the proposition.
\end{proof}

The coupling has the particular independence needed in
Section~\ref{sec:stability}: the possible edit locations can be sampled
before the pure spin field is examined.  Conditional refreshed values
remain unrestricted.  Thus a later estimate that permits arbitrary spin
assignments in the marked squares applies to this exact perturbed law,
up to the vanishing exceptional probability already quantified above.
